\documentclass[10pt,twocolumn]{article}
\usepackage[T1]{fontenc}
\usepackage[utf8]{inputenc}
\usepackage{lmodern}
\usepackage[margin=1.8cm,top=2.4cm,bottom=2cm,columnsep=0.6cm]{geometry}
\usepackage{fancyhdr}
\usepackage{titlesec}
\usepackage{booktabs}
\usepackage{amsmath,amssymb,amsfonts}
\usepackage{microtype}
\usepackage{xcolor}
\usepackage{mdframed}
\usepackage{parskip}
\usepackage{enumitem}
\usepackage{hyperref}

\definecolor{rulered}{RGB}{180,30,30}
\definecolor{boxgray}{RGB}{245,245,245}
\definecolor{keywordgray}{RGB}{100,100,100}

\pagestyle{fancy}
\fancyhf{}
\fancyhead[L]{\textsc{\small Claude Mag}}
\fancyhead[C]{\small\textit{Machine Learning Research Digest}}
\fancyhead[R]{\small 2026-08-19}
\fancyfoot[C]{\small\thepage}
\renewcommand{\headrulewidth}{0.8pt}

\titleformat{\section}{\normalsize\bfseries\color{rulered}}{}{0pt}{}%
  [\vspace{-4pt}\color{rulered}\rule{\columnwidth}{0.4pt}]
\titlespacing{\section}{0pt}{8pt}{4pt}

\hypersetup{colorlinks=true,urlcolor=rulered,linkcolor=black}

\begin{document}
\setlength{\parindent}{0pt}
\setlength{\parskip}{4pt}

{\small\color{keywordgray}\textbf{Keywords:}
  efficient inference \textbullet{}
  KV-cache compression \textbullet{}
  process reward model \textbullet{}
  chain-of-thought reasoning}\par\vspace{4pt}

{\LARGE\bfseries\raggedright
  Fewer Tokens, Smaller Cache:
  Reward-Coordinated Efficient Reasoning}\par\vspace{4pt}

{\large\itshape\raggedright
  Process Reward Signals Jointly Guide Cache Eviction and Generation
  Length to Halve Inference Cost Without Sacrificing Accuracy}\par\vspace{6pt}

{\small\textbf{By} Qiyuan Zhu, Dezhi Li, Pengyu Cheng, Tianle Chen,
  Jiacheng Wang, Ruijie Shen, Hao Gu, Sida Lin, Zirui Liu, Jiacheng Liu,
  Sirui Han}\par\vspace{2pt}

{\small\color{keywordgray}arXiv:2608.04771 \textbullet{} Preprint, August 2026}
\par\vspace{2pt}\noindent{\color{rulered}\rule{\columnwidth}{0.6pt}}\vspace{6pt}

Large Reasoning Models generate long chain-of-thought traces that strain
inference budgets in two entangled ways: the KV-cache grows with context length,
and token generation itself consumes compute. Existing compression methods attack
only the cache and apply uniform eviction policies across the entire trajectory.
This paper shows that a process reward model (PRM) naturally tracks
\emph{where} compression is safe and \emph{when} generation can stop,
enabling simultaneous reduction of both costs under a single reward-coordinated
framework called RCER.

\begin{mdframed}[backgroundcolor=boxgray,linecolor=rulered,linewidth=1pt,
    innerleftmargin=6pt,innerrightmargin=6pt,
    innertopmargin=6pt,innerbottommargin=6pt]
\textbf{\small AT A GLANCE}\par\vspace{4pt}
\begin{description}[leftmargin=0pt,labelsep=4pt,itemsep=2pt,parsep=0pt,
    font=\small\bfseries]
  \item[Problem:] LRM inference is expensive in both KV-cache memory and
    generated tokens; uniform compression sacrifices accuracy.
  \item[Why it matters:] Efficient LRM deployment is blocked by inference cost;
    a principled two-sided solution would substantially lower the barrier.
  \item[Method:] Use a process reward model to guide KV-cache eviction
    (high-reward steps tolerate more deletion) and generation stopping
    (stop when the reward signals convergence).
  \item[Result:] Simultaneous reduction in cache size and token count,
    with accuracy preserved at compression rates that defeat uniform baselines.
  \item[Evidence:] Ablations confirm that reward-guided eviction outperforms
    attention-score and recency heuristics at matched compression ratios.
\end{description}
\end{mdframed}

\section{The Big Picture}

Large Reasoning Models (LRMs) are chain-of-thought systems trained to produce
extended internal deliberation before outputting a final answer. The deliberation
is effective — LRMs achieve state-of-the-art performance on mathematical,
scientific, and code reasoning tasks — but it is expensive. A single
difficult-problem response may span thousands of tokens of intermediate reasoning,
all of which must be stored in the KV-cache during generation and computed token
by token.

Two inference costs accumulate jointly: the KV-cache grows linearly with the
generated sequence, consuming memory bandwidth and limiting batch size; and the
number of generated tokens determines total compute via floating-point operations.
These costs are interdependent: a smaller cache means the model has less context,
which causes it to regenerate lost information and produce more tokens than it
would with a full cache.

This paper identifies the joint nature of the problem and proposes RCER
(Reward-Coordinated Efficient Reasoning) as the first method to simultaneously
minimize both quantities under a single process reward signal.

\section{Why Existing Approaches Fall Short}

\textbf{Cache-only compression} reduces KV-cache size by evicting past tokens
based on heuristics: recency (evict oldest tokens), attention score (evict
least-attended tokens), or importance scores derived from gradient norms.
These approaches are unaware of the reasoning trajectory's progress and evict
uniformly, even at critical steps where context loss causes the model to
misidentify its reasoning state.

\textbf{Generation length reduction} without cache-awareness typically operates
through early-exit mechanisms that terminate generation when some threshold is
reached. Without understanding the quality of the reasoning state — only its
length — these mechanisms can terminate prematurely, before the model has reached
a confident conclusion.

Neither class of method accounts for the \emph{interaction}: compressing the
cache causes more generation, and shortening generation can leave the cache
unnecessarily large. RCER resolves this by tying both decisions to the same
quality signal.

\section{Core Mechanism}

RCER augments standard LRM decoding with a process reward model $\phi$
that is evaluated at each reasoning step $k$ to produce a scalar quality
estimate $r_k = \phi(x, t_{\leq k}) \in [0, 1]$, where $x$ is the input
problem and $t_{\leq k}$ is the reasoning trace up to step $k$.

The PRM drives two decisions:
\begin{enumerate}[itemsep=2pt,parsep=0pt]
  \item \textbf{Cache eviction:} When $r_k$ is high (the model has committed to
    a correct sub-conclusion), the cached representations of older tokens are
    less necessary — the model's trajectory is converging. RCER evicts an
    amount proportional to $r_k$ from the oldest entries in the cache.
  \item \textbf{Generation stopping:} When $r_k$ exceeds a convergence threshold
    $\tau$ and has been stable for $\delta$ steps, RCER terminates generation
    and commits the current answer, avoiding wasteful continuation of a
    completed reasoning chain.
\end{enumerate}

No additional model beyond the PRM (which may be a lightweight verifier) is
required.

\section{Technical Formulation}

Let $\pi_\theta$ be the LRM and $r_k = \phi(x, t_{\leq k})$ the PRM score at step $k$.

\textbf{Cache eviction budget at step $k$:}
\[
  b_k = \beta \cdot r_k \cdot |\mathcal{C}_k|
\]
where $|\mathcal{C}_k|$ is the current cache size and $\beta$ is a
compression rate hyperparameter. Tokens are evicted from the oldest
$b_k$ positions.

\textbf{Generation stopping criterion:}
\[
  \text{stop if} \quad r_k \geq \tau \quad \text{and} \quad
  \frac{1}{\delta}\sum_{j=k-\delta}^{k} r_j \geq \tau - \varepsilon
\]
This ensures both that the current state is high-quality and that the
quality has been stable (not just a local spike).

\textbf{Effective inference cost:}
\[
  \mathcal{L} = \underbrace{\sum_{k=1}^{K} |\mathcal{C}_k|}_{\text{cache cost}}
  + \lambda \cdot \underbrace{K}_{\text{token cost}}
\]
RCER minimizes $\mathcal{L}$ under the constraint that accuracy remains
within $\epsilon$ of unconstrained decoding.

\section{Learning or Inference Procedure}

RCER is a \emph{training-free} inference-time method. The LRM $\pi_\theta$
and PRM $\phi$ are both used off-the-shelf without any fine-tuning. The
only hyperparameters are:
\begin{itemize}[itemsep=2pt,parsep=0pt]
  \item $\beta \in [0, 1]$: eviction rate (ablated: 0.1, 0.3, 0.5)
  \item $\tau \in [0, 1]$: convergence threshold (ablated: 0.7, 0.8, 0.9)
  \item $\delta$: stability window in steps (ablated: 5, 10, 20)
  \item $\lambda$: weight trading off cache vs.\ token cost
\end{itemize}

\section{What the Guarantee Says}

No formal guarantee is provided. The empirical claim is that RCER
simultaneously reduces both KV-cache occupancy and generated tokens
while preserving accuracy within a specified tolerance of unconstrained
decoding. The main theoretical insight — that PRM score correlates with
the reasoning state's tolerance to context loss — is supported by
ablation experiments rather than formal analysis.

\section{Experimental Findings}

\begin{itemize}[itemsep=2pt,parsep=0pt]
  \item \textbf{Cache reduction:} RCER reduces peak KV-cache size by
    40--55\% on GSM8K, MATH-500, and LiveCodeBench at accuracy
    within 1\% of unconstrained decoding.
  \item \textbf{Token reduction:} RCER simultaneously reduces generated
    tokens by 20--35\% through early stopping at convergence.
  \item \textbf{Uniform baselines fail:} At 40\% cache compression,
    attention-score eviction drops accuracy by 4--8\%; RCER drops by less
    than 1\%.
  \item \textbf{Interaction confirmed:} Experiments show that cache-only
    compression at rate $\beta$ increases token count by 10--25\%,
    confirming the interaction and the necessity of joint coordination.
\end{itemize}

\section{Ablations and Interpretation}

\textbf{Eviction strategy:} Reward-guided eviction consistently outperforms
recency, attention-score, and random eviction at matched budgets.
The advantage grows with compression ratio, indicating that the reward
signal is most critical at high compression where wrong eviction choices
compound.

\textbf{PRM quality:} A lightweight 1B-parameter verifier achieves 85--90\%
of the benefit of a full-scale PRM, suggesting the approach is practical
even without access to a high-accuracy reward model.

\textbf{Stopping ablation:} Generation stopping without cache coordination
saves tokens but wastes cache — the combined approach is strictly better
than either alone on the joint cost metric.

\section*{Reference}

Qiyuan Zhu, Dezhi Li, Pengyu Cheng, Tianle Chen, Jiacheng Wang, Ruijie Shen,
Hao Gu, Sida Lin, Zirui Liu, Jiacheng Liu, Sirui Han.
\textbf{Fewer Tokens, Smaller Cache: Reward-Coordinated Efficient Reasoning}.
arXiv:2608.04771, August 2026.
\url{https://arxiv.org/abs/2608.04771}

\end{document}
